\subsection{拉普拉斯反变换}
\label{subsec:inverse_laplace_transform}
拉普拉斯反变换(亦称逆变换)写为
\begin{equation}
    f(t) = \mathcal{L}^{-1} [ \bar{f}(p)].
\end{equation}
由傅里叶逆变换导出的具体表达式是
\begin{equation}
    f(t)=\frac{1}{2 \pi \imath} \int_{\sigma-\imath \infty}^{\sigma+\imath \infty} \bar{f}(p) e^{ p t} dp,
    \label{eq:bromwich_repeat}
\end{equation}
见式\eqref{eq:fourier_mellin_integral}.该式是著名的Bromwich积分,也叫做\textbf{Riemann--Mellin积分}.

积分路径是$p$平面上的竖直直线$\Re p=\sigma$,方向从$\sigma-\imath\infty$到$\sigma+\imath\infty$.
须取$\sigma>\sigma_0$,使$\bar{f}$的全部奇点落在该线左侧;右半平面$\Re p>\sigma_0$上$\bar{f}$解析.
$t>0$时用左侧大圆弧把Bromwich线闭上,在约当引理条件下圆弧贡献趋于零,原函数等于闭路内留数之和.

\textbf{推广的约当引理.}
设 $C_R$ 是以 $p=0$ 为圆心、半径为 $R$ 的圆周在直线 $\Re p=a(>0)$ 左侧的圆弧.
若当 $|p| \rightarrow \infty$ 时, $\bar{f}(p)$ 在
$\frac{\pi}{2}-\delta \leqslant \operatorname{Arg} p \leqslant \frac{3}{2} \pi+\delta$
中一致趋于零($\delta$ 是小于 $\frac{\pi}{2}$ 的任意正数),则
$$
\lim _{R \rightarrow \infty} \int_{C_R} \bar{f}(p) e^{p t} d p=0 \quad(t>0) .
$$
于是对$t>0$,
\begin{equation}
    f(t) = \sum_{j}\operatorname{Res}\bigl[\bar{f}(p)\,e^{pt},\,p_j\bigr],
\end{equation}
其中求和取闭路所围的全部奇点.简单极点$p_j$上
\[
  \operatorname{Res}
  =\lim_{p\to p_j}(p-p_j)\,\bar{f}(p)\,e^{pt};
\]
$n$阶极点则对$(p-p_j)^n\bar{f}(p)e^{pt}$求$n-1$阶导后再除以$(n-1)!$.

\begin{example}
利用Bromwich积分求$1/\sqrt{p}$的原函数.
\end{example}
\begin{solution}
被积函数$e^{pt}/\sqrt{p}$的唯一奇点是$p=0$,这是支点而非极点.
不能直接用左侧半圆闭上Bromwich线.
先自$p=0$沿负实轴把平面切开,取辐角主值$\operatorname{Arg}p\in(-\pi,\pi)$,
再作回路:上行Bromwich线$\to$左侧大圆弧$C_R\to$割线上沿$l_2\to$绕原点的小圆$C_\varepsilon\to$割线下沿$l_1$.
被积函数在该单值分支上于回路内部解析,故回路积分为零.

令$R\to\infty$,由推广的约当引理$\int_{C_R}\to 0$.
小圆上弧长$2\pi\varepsilon$,$|e^{pt}/\sqrt{p}|$量级为$\varepsilon^{-1/2}$,故$\varepsilon\to 0$时$\int_{C_\varepsilon}\to 0$.
因此Bromwich积分等于割线两侧贡献的相反数:
\begin{equation}
f(t)=-\frac{1}{2 \pi \imath} \lim _{R \rightarrow \infty}\left(\int_{l_1}+\int_{l_2}\right) \frac{e^{p t}}{\sqrt{p}} d p .
\end{equation}
割线下沿辐角为$-\pi$,上沿为$\pi$,于是
$\sqrt{p}=-\imath\sqrt{|p|}$(下沿),$\sqrt{p}=\imath\sqrt{|p|}$(上沿).代入后两项同号,
\begin{equation}
\begin{aligned}
f(t)
& =\frac{1}{\pi} \int_{-\infty}^0 e^{p t} \frac{1}{\sqrt{|p|}} d p .
\end{aligned}
\end{equation}
令$y=\sqrt{|p|\,t}$,则$p=-y^2/t$,$\mathrm{d}p=-(2y/t)\,\mathrm{d}y$,
\begin{equation}
\begin{aligned}
f(t)
& =\frac{1}{\pi} \int_{\infty}^0 e^{-y^2} \frac{\sqrt{t}}{y}\left(-\frac{2 y}{t} d y\right)
 =\frac{2}{\pi \sqrt{t}} \int_0^{\infty} e^{-y^2} d y
 =\frac{1}{\sqrt{\pi t}} .
\end{aligned}
\end{equation}
亦即$\mathcal{L}\{1/\sqrt{\pi t}\}=1/\sqrt{p}$.
\end{solution}

很多时候不必直接做围道.若$f(t)$满足存在条件,则可在积分号下对$p$求导:
\begin{equation}
\bar{f}^{(n)}(p)=\int_0^{\infty}(-t)^n f(t) e^{-p t} d t,
\end{equation}
所以
\begin{equation}
    \bar{f}^{(n)}(p)\risingdotseq (-t)^n f(t).
\end{equation}
特别地由$\mathcal{L}\{1\}=1/p$得
\[
    \begin{aligned}
    & \frac{1}{p^2}=-\frac{d}{d p} \frac{1}{p} \risingdotseq  t, \\
    & \frac{1}{p^3}=\frac{1}{2} \frac{d^2}{d p^2} \frac{1}{p} \risingdotseq \frac{1}{2} t^2 .
    \end{aligned}
\]

有理像函数可先裂项再逐项反演.例如
\begin{equation}
    \begin{aligned}
    \frac{1}{p^3(p+\alpha)}
    & =\frac{1}{\alpha} \frac{1}{p^3}-\frac{1}{\alpha^2} \frac{1}{p^2}+\frac{1}{\alpha^3} \frac{1}{p}-\frac{1}{\alpha^3} \frac{1}{p+\alpha} \\
    & \risingdotseq \frac{t^2}{2 \alpha}-\frac{t}{\alpha^2}+\frac{1}{\alpha^3}-\frac{e^{-\alpha t}}{\alpha^3}.
    \end{aligned}
\end{equation}

\begin{example}
求 $\dfrac{e^{-\alpha p}}{p(p+b)}$ 的原函数($\alpha>0$,$b\neq 0$).
\end{example}
\begin{solution}
$\dfrac{1}{p} \risingdotseq \Theta(t)$.应用延迟定理,
\[
\frac{e^{-\alpha p}}{p} \risingdotseq \Theta(t-\alpha),
\qquad
\frac{1}{p+b} \risingdotseq e^{-b t}.
\]
像函数乘积对应卷积,
\begin{equation}
\begin{aligned}
\frac{e^{-\alpha p}}{p(p+b)}
& \risingdotseq \int_0^t \Theta(\tau-\alpha) e^{-b(t-\tau)} d \tau \\
& =\Theta(t-\alpha) \int_\alpha^t e^{-b(t-\tau)} d \tau \\
& =\frac{1-e^{-b(t-\alpha)}}{b}\,\Theta(t-\alpha).
\end{aligned}
\end{equation}
也可先把分母裂成$\frac{1}{b}\bigl(\frac{1}{p}-\frac{1}{p+b}\bigr)$,再对每一项用延迟定理核对.
\end{solution}
